\chapter{Soundness Is a Fold}\label{ch:sound}

A rewrite rule that terminates but changes the value is worse than useless: it is confident and
wrong. This chapter is about the theorems that say a rule \emph{means} something --- first over a
small fragment where everything is decidable, then over the real numbers --- and about the two
rules that turned out not to, whose refutations are the most instructive theorems in the engine.

\section{Soundness for a rule, and for a run}

Fix a relation \lc{rel} between terms that we intend to mean ``these denote the same thing''. A
rule is sound for \lc{rel} if its output is related to its input. That is a statement about one
firing at one node. What the notebook needs is a statement about the \emph{whole run}: the input
of a cell is related to its output. The passage from one to the other is a fold over the
rewriter's recursion, and it needs exactly four properties of the relation.

\begin{minted}{lean}
structure Congruence where
  rel : Expr → Expr → Prop
  refl : ∀ e, rel e e
  trans : ∀ {a b c}, rel a b → rel b c → rel a c
  /-- Rewriting each child soundly rewrites the node soundly. -/
  congr : ∀ e cs, RelList rel (children e) cs → rel e (withChildren e cs)
  /-- Reordering a sum's or product's arguments is sound — what `canon` does silently. -/
  canon : ∀ e, rel e (canon e)

def RuleSoundFor (C : Congruence) (r : Rule W) : Prop :=
  ∀ e res, r.apply e = some res → C.rel e res.result

theorem normalize_sound_for (C : Congruence) (rules : List (Rule W))
    (hs : ∀ r ∈ rules, RuleSoundFor C r) (e : Expr) (s : Array Step) :
    C.rel e ((normalize rules e).run' s)
\end{minted}

The theorem is proved once, for an abstract \lc{Congruence}, by following the rewriter's
recursion: children are rewritten (congruence), the node is put into canonical order
(\lc{canon}), a rule fires (its own soundness), the result is rewritten again (transitivity,
and the recursion). Nothing in the proof knows what \lc{rel} is. That was a deliberate
generalization made in the third week, and it paid three times: the integers, the reals, and
the reals under a derivative each instantiate the same fold.

\section{The integer fragment}

The first semantics is the smallest one that can be \emph{decided}: an evaluator
\lc{eval?}, from an environment and a term to an optional integer, that gives a value to numerals with denominator one,
variables from an environment, sums, products and natural powers, and \lc{none} to everything
else --- functions, fractions, matrices. Two terms are related when the second refines the first:
wherever the first has a value, the second has the same one.

\begin{minted}{lean}
def Refines (e e' : Expr) : Prop := ∀ ρ v, eval? ρ e = some v → eval? ρ e' = some v
theorem simplify_sound (e : Expr) : Refines e (simplify0 e)
\end{minted}

Every one of the seven simplification rules is sound for this relation, and the fold gives
\lc{simplify_sound}. On its own this is a modest claim --- it says nothing about $\sqrt 2$ or
$\sin$ --- but it was the first end-to-end theorem about the engine, checked on the first day the
rules existed, and it fixed the \emph{shape} every later semantics would take.

\section{The reals}

The real semantics lives in \file{proofs/}, because $\mathbb{R}$ is Mathlib's. It is a total
function: powers are \lc{Real.rpow}, \texttt{sqrt} is \lc{Real.sqrt}, \texttt{ln} is
\lc{Real.log}, and where these are undefined in mathematics they take Mathlib's \emph{junk} values
($\log 0 = 0$, $\sqrt{-1} = 0$, $0^{-1} = 0$). Junk values are not a weakness here; they are what
makes every statement below an unconditional equation between total functions, so that the side
conditions a rule really needs have to be \emph{written down} rather than hidden inside a partial
function.

\begin{minted}{lean}
def evalR (ρ : EnvR) : Expr → ℝ
  | .num q => (q.val : ℝ)
  | .var x => ρ x
  | .add es => sumR ρ es
  | .mul es => prodR ρ es
  | .pow b e => (evalR ρ b) ^ (evalR ρ e)
  | .fn f [a] => applyFn f (evalR ρ a)
  | .fn _ _ => 0
  | .matrix _ => 0

def SemEqR (e e' : Expr) : Prop := ∀ ρ, evalR ρ e = evalR ρ e'
\end{minted}

A bridge theorem, \lc{evalR_of_eval?}, says that wherever the integer evaluator is defined the two
agree; so the theorem of the previous section was a restriction of the truth, not a different
claim. The reader will meet this pattern again in Chapter~\ref{ch:deriv}: three semantics, each
shown to be a restriction of the next.

\section{Two rules that are not sound}

Five of the seven rules are sound for \lc{SemEqR} and the fold applies to them. Two are not, and
the interesting part is that this is \emph{proved}.

\begin{thmbox}[\lc{not_collectPowers_soundR}]
  The rule $b^m \cdot b^n \to b^{m+n}$ is not unconditionally sound over $\mathbb{R}$: at $x = 0$
  it rewrites $x \cdot x^{-1}$, which is $0$, to $x^0$, which is $1$.
\end{thmbox}

\begin{minted}{lean}
theorem not_collectPowers_soundR : ¬ RuleSoundR collectPowers := by
  intro hs
  have h := hs (.mul [.var "x", .pow (.var "x") (.num (Q.ofInt (-1)))]) _ rfl (fun _ => 0)
  norm_num [evalR_mul, prodR_cons, prodR_nil, evalR_var, evalR_pow, evalR_num, …] at h
\end{minted}

\begin{thmbox}[\lc{not_functionRules_soundR}]
  The rule $\exp(\ln x) \to x$ is not unconditionally sound: at $x = -1$ the left side is
  $\exp(\ln(-1)) = \exp(\ln 1) = 1$, because \lc{Real.log} is even, while the right side is $-1$.
\end{thmbox}

Both rules are the standard convention of every computer-algebra system, and the engine keeps
them. What changed is the label. The collect-powers rule is \status{conditional}, and the theorem
\lc{collectPowers_soundR_on} says exactly when it is right: whenever the merged base is positive.
The notebook shows that condition next to the step. A reader who has only ever seen $x/x = 1$
printed without comment now sees the $x \neq 0$ that was always there.

\begin{logbox}[On pinning Mathlib]
  The proofs package pins Mathlib at the release tag whose \texttt{lean-toolchain} equals the
  engine's. The first plan had been ``pin to Mathlib master''; on the day, master was on a
  release candidate \emph{ahead} of stable, so following it would have forced the engine off the
  policy of latest stable. Pin to the tag that matches; bump both together.
\end{logbox}

\section{What the fold does not cover}

The fold is stated for rules that are sound \emph{unconditionally}. A conditional rule cannot be
folded, because the hypothesis at one step is about the values at that step, and the fold does
not know them. So the notebook's ``this derivation is sound'' is not one theorem; it is the
conjunction of the statuses of the steps, exactly as displayed. Chapter~\ref{ch:integrate} turns
this into a precise statement --- the integral checker's theorem takes ``the derivation shown was
sound at the point'' as its hypothesis, and that hypothesis is what the statuses say.

\begin{trybox}
\begin{minted}{text}
x * x^(-1)
exp(ln(x))
\end{minted}
  Both simplify. Click each step and read the proof status: \status{conditional}, with the
  hypothesis written out. Then try \texttt{subst(x * x\^{}(-1), x, 0)}: the answer is $1$,
  because the children of a command are normalized first and the conditional step has already
  fired. The label on that step is the only warning, which is the point of the label.
\end{trybox}

\section{Exercises}

\begin{exercisebox}[A third refutation]
  The rule $\sqrt{x^2} \to x$ is not in the engine. Prove in the style of the two refutations
  above that it would not be unconditionally sound over $\mathbb{R}$, and state the side
  condition under which it is.
\end{exercisebox}

\begin{exercisebox}[Congruence laws]
  Show that \lc{SemEqR} satisfies the \lc{canon} law: a permutation of the arguments of a sum
  does not change \lc{sumR}. (The engine's proof is \lc{sumR_perm}, an induction on
  \lc{List.Perm}.) Which of the four laws would fail for a relation that compared terms only at
  the single point $\rho = 0$?
\end{exercisebox}
